\chapter{Modelisation mathematique du risque}
\label{chap:modelisation}

Ce chapitre formalise le modele de risque implemente par
\texttt{RiskCalculationService} et \texttt{DashboardMetricsService}. Chaque formule
est d'abord posee, puis accompagnee de proprietes demontrees qui justifient les
choix de conception (bornage, monotonie, coherence de la classification) et qui
mettent en evidence, lorsque c'est le cas, les ecarts assumes entre la specification
initiale (\emph{pfa\_formule.pdf}) et son implementation.

\section{Notations}

\begin{longtable}{@{}p{2.6cm}p{10.2cm}@{}}
\toprule
\textbf{Symbole} & \textbf{Signification} \\
\midrule
\endhead
$h$ & un actif (hote) reconcilie entre Qualys et la CMDB \\
$ACS(h)$ & Asset Criticality Score de $h$, $ACS(h)\in[1,5]$ \\
$i \in \{c,h,m,l\}$ & indice de severite : Critique, Elevee (\emph{High}), Moyenne,
Faible \\
$w_i$ & poids configure de la severite $i$ \\
$QDS$ & Qualys Detection Score d'une vulnerabilite, $QDS\in[0,100]$ \\
$\mathrm{Avg}(QDS_i)$ & moyenne des $QDS$ des vulnerabilites ouvertes de severite $i$
sur $h$ \\
$n_i$ & nombre de vulnerabilites de severite $i$ sur $h$ \\
$I_{External}$ & facteur d'exposition reseau (Internet) \\
$M_{SOC}$, $M_{AV}$, $M_{Patch}$ & facteurs de modulation lies aux controles
compensatoires \\
$ARS(h)$ & Asset Risk Score de $h$ \\
$N$ & nombre total d'actifs actifs et scannes \\
$TTR$ & temps de resolution (\emph{Time To Remediate}) d'une vulnerabilite corrigee,
en jours \\
$\mathrm{MTTR}_w$ & temps moyen de remediation pondere par classe de risque \\
$ERS$ & Enterprise Risk Score \\
\bottomrule
\end{longtable}

\section{Le score de criticite d'actif (ACS)}

\subsection{Definition}

L'ACS traduit l'impact metier d'un actif independamment de ses vulnerabilites : un
routeur de coeur de reseau en production n'a pas la meme importance qu'un serveur de
test. Il est deduit de deux poids configures par l'administrateur -- un poids par
role technique $\rho: \mathrm{Role}\to[0,1]$ et un poids par environnement
$\varepsilon: \mathrm{Env}\to[0,1]$ -- moyennes puis recalees sur l'intervalle
$[1,5]$ :

\begin{equation}
ACS(h) = 1 + 4 \cdot \underbrace{\frac{\rho(\mathrm{role}(h)) + \varepsilon(\mathrm{env}(h))}{2}}_{\displaystyle \bar{x}(h)\,\in\,[0,1]}
\label{eq:acs}
\end{equation}

\begin{propriete}[Unicite du recalage affine]
\label{prop:acs-affine}
La fonction $f(x) = 1+4x$ est l'unique fonction affine $f(x)=ax+b$ verifiant
$f(0)=1$ et $f(1)=5$, et elle envoie bijectivement $[0,1]$ sur $[1,5]$ en preservant
l'ordre.
\end{propriete}

\begin{proof}
Une fonction affine est entierement determinee par ses valeurs en deux points
distincts. Les conditions aux bornes imposent $f(0)=b=1$, puis
$f(1)=a\cdot 1+b=5 \Rightarrow a=4$ : le couple $(a,b)=(4,1)$ est donc l'unique
solution, d'ou $f(x)=1+4x$. Comme $a=4>0$, $f$ est strictement croissante, donc
injective sur $[0,1]$ ; combinee a $f(0)=1$ et $f(1)=5$, elle realise une bijection
croissante de $[0,1]$ sur $[1,5]$. En particulier, pour tous $x,x'\in[0,1]$,
$x\le x' \iff f(x)\le f(x')$ : l'ordre des criticites relatives entre actifs est
exactement preserve par le recalage.
\end{proof}

Cette propriete justifie le choix d'implementation : comparer deux $ACS$ revient
exactement a comparer les moyennes $\bar x$ sous-jacentes, aucune information
d'ordre n'est perdue ni deformee par le recalage.

\section{Le score de vulnerabilite pondere par severite}

\subsection{Formule de reference}

La specification (eq. 1, terme interieur) pondere la moyenne des $QDS$ de chaque
classe de severite par un poids fixe :

\begin{equation}
\mathrm{Score}_{\text{ref}}(h) = \sum_{i\in\{c,h,m,l\}} w_i \cdot \mathrm{Avg}(QDS_i)
\label{eq:qds-ref}
\end{equation}

avec, valeurs de reference, $w_c=1.0$, $w_h=0.2$, $w_m=0.05$, $w_l=0.01$.

\subsection{Raffinement implemente}

L'implementation (\texttt{RiskCalculationService.computeWeightedQds}) introduit un
terme correctif lie au \emph{nombre} de vulnerabilites de chaque severite, afin
qu'une classe comportant davantage de failles ouvertes pese un peu plus lourd, sans
pour autant laisser une severite anecdotique (une seule faille isolee) ou au
contraire un tres grand nombre de failles de meme severite faire exploser le score de
maniere disproportionnee :

\begin{equation}
\mathrm{WeightedQDS}(h) = \sum_{i\in\{c,h,m,l\}} w_i \cdot \mathrm{Avg}(QDS_i) \cdot n_i^{1/100}
\label{eq:qds-impl}
\end{equation}
\noindent (convention : le terme vaut $0$ si $n_i = 0$.)

\begin{propriete}[Le facteur de comptage est croissant mais quasi borne en pratique]
\label{prop:count-damping}
Pour $n\ge 1$, la fonction $\varphi(n)=n^{1/100}$ est strictement croissante et
verifie $\varphi(1)=1$. De plus, pour tout $n\le 10^{4}$ (soit dix mille
vulnerabilites d'une meme severite sur un meme hote, un majorant tres large en
pratique), $\varphi(n) \le 1{,}0965$.
\end{propriete}

\begin{proof}
$\varphi$ est derivable sur $(0,\infty)$ avec
$\varphi'(n) = \tfrac{1}{100}n^{-99/100} > 0$, donc $\varphi$ est strictement
croissante ; $\varphi(1)=1^{1/100}=1$.
Etant croissante, $\varphi(n)\le\varphi(10^4)$ pour tout $n\le 10^4$. Or
\[
\varphi(10^4) = (10^4)^{1/100} = 10^{4/100} = 10^{0{,}04}
             = e^{0{,}04\ln 10} = e^{0{,}09210\ldots} \approx 1{,}0965 .
\]
\end{proof}

\begin{corollaire}[Le raffinement domine strictement un multiplicateur brut]
\label{cor:count-vs-raw}
Remplacer $n_i^{1/100}$ par un multiplicateur brut $n_i$ (comme le suggererait une
ponderation « poids $\times$ moyenne $\times$ nombre ») fait croitre le terme
lineairement et sans borne avec $n_i$, alors que \cref{prop:count-damping} montre que
$n_i^{1/100}$ reste inferieur a $1{,}1$ meme pour $n_i$ de l'ordre du millier. Par
exemple, pour $n_i=1000$~: le multiplicateur brut vaut $1000$ (le terme est
multiplie par mille) contre $1000^{1/100}=10^{0{,}03}\approx 1{,}0715$
(augmentation de $7\%$ seulement) avec le raffinement implemente.
\end{corollaire}

Ce corollaire justifie formellement le commentaire de conception du code source
(\emph{« gives classes with more open vulnerabilities a slowly growing extra push,
without letting a single stray finding swing the term »}) : le nombre d'occurrences
influence le score de maniere monotone (\cref{prop:count-damping}), donc coherente
avec l'intuition « plus de failles ouvertes = plus de risque », mais son effet est
strictement borne en pratique (\cref{cor:count-vs-raw}), de sorte que ni une faille
isolee ni une accumulation massive de failles de faible impact ne peuvent, a elles
seules, piloter le score.

\section{L'Asset Risk Score (ARS)}

\subsection{Equation fondamentale}

\begin{equation}
ARS(h) = \min\Big(\; ACS(h)\times \mathrm{WeightedQDS}(h) \times I_{External}(h) \times M_{Controls}(h)\;,\;1000 \Big)
\label{eq:ars}
\end{equation}

avec le facteur de controles compensatoires :

\begin{equation}
M_{Controls}(h) = M_{SOC}(h)\times M_{AV}(h)\times M_{Patch}(h)
\label{eq:mcontrols}
\end{equation}

Chaque facteur de modulation est applique de maniere \emph{conditionnelle}, avec
$1$ (element neutre de la multiplication) comme valeur par defaut lorsque sa
condition d'activation n'est pas remplie :

\begin{align}
I_{External}(h) &= \begin{cases} 1{,}2 & \text{si $h$ expose sur Internet (QID 45004)} \\ 1{,}0 & \text{sinon} \end{cases} \\
M_{SOC}(h) &= \begin{cases} 0{,}8 & \text{si supervise SOC} \\ 1{,}0 & \text{sinon} \end{cases} \\
M_{AV}(h) &= \begin{cases} 0{,}9 & \text{si EDR/antivirus present} \\ 1{,}0 & \text{sinon} \end{cases} \\
M_{Patch}(h) &= \begin{cases} 1{,}3 & \text{si exclu du cycle de patch standard} \\ 1{,}0 & \text{sinon} \end{cases}
\end{align}

Le \cref{tab:multiplicateurs} recapitule ces valeurs de reference (table 1 de la
specification).

\begin{table}[H]
\centering
\caption{Matrice des multiplicateurs de risque residuel}
\label{tab:multiplicateurs}
\begin{tabular}{@{}lllr@{}}
\toprule
\textbf{Facteur} & \textbf{Source} & \textbf{Condition} & \textbf{Valeur} \\
\midrule
$I_{External}$ & Qualys (QID 45004) & si expose sur Internet & $1{,}2$ \\
$M_{SOC}$ & CMDB & si actif monitore (SIEM) & $0{,}8$ \\
$M_{AV}$ & CMDB & si presence d'un EDR/antivirus & $0{,}9$ \\
$M_{Patch}$ & CMDB & si exclusion du cycle de patch & $1{,}3$ \\
\bottomrule
\end{tabular}
\end{table}

\subsection{Positivite et monotonie}

\begin{propriete}[Positivite]
\label{prop:ars-pos}
Pour tout $h$, $ARS(h) \in [0,1000]$.
\end{propriete}

\begin{proof}
$ACS(h)\ge 1>0$ par \cref{prop:acs-affine}. $\mathrm{WeightedQDS}(h)$ est une somme de
termes $w_i\cdot\mathrm{Avg}(QDS_i)\cdot n_i^{1/100}$ tous positifs ou nuls (poids
$w_i\ge0$, $QDS\ge0$, et $n_i^{1/100}\ge0$ ou terme conventionnellement nul si
$n_i=0$), donc $\mathrm{WeightedQDS}(h)\ge0$. $I_{External}(h)$ et $M_{Controls}(h)$
sont des produits de constantes strictement positives, donc strictement positifs.
Le produit $ACS(h)\times\mathrm{WeightedQDS}(h)\times I_{External}(h)\times
M_{Controls}(h)$ est donc positif ou nul, et $\min(\cdot,1000)$ le majore par $1000$.
\end{proof}

\begin{propriete}[Monotonie]
\label{prop:ars-mono}
$ARS(h)$ est une fonction (faiblement) croissante de chacun de ses quatre facteurs
$ACS$, $\mathrm{WeightedQDS}$, $I_{External}$, $M_{Controls}$, les trois autres etant
fixes.
\end{propriete}

\begin{proof}
Notons $P = ACS\times \mathrm{WeightedQDS}\times I_{External}\times M_{Controls}$ le
produit avant saturation, et $g(x)=\min(x,1000)$. Soient $a\le a'$ deux valeurs
possibles pour l'un des quatre facteurs (les trois autres, notes $k\ge0$, etant
fixes) : alors $P(a) = k\cdot a \le k\cdot a' = P(a')$ car $k\ge0$. La fonction
$g$ est elle-meme croissante : pour $x\le x'$, $\min(x,1000)\le\min(x',1000)$ (cas
immediat par disjonction $x'\le1000$ ou $x'>1000$). La composee d'une fonction
croissante d'une variable ($P(\cdot)$ dans la direction consideree) avec une fonction
croissante ($g$) est croissante, donc $ARS = g\circ P$ est croissante en chacun de
ses facteurs.
\end{proof}

Cette propriete garantit que le modele ne peut jamais produire de contre-intuition
grave : ajouter une mesure de protection (facteur $<1$), reduire l'exposition
($I_{External}=1$ au lieu de $1{,}2$), ou reduire la gravite des vulnerabilites ne
peut jamais \emph{augmenter} l'ARS affiche.

\subsection{Condition de saturation}

Le plafond a $1000$ evite qu'un score sans borne ne perde toute signification
au-dela du pire cas possible. On peut caracteriser precisement quand il s'active.

\begin{propriete}[Seuil de saturation]
\label{prop:saturation}
Avec les valeurs de reference du \cref{tab:multiplicateurs} et $ACS(h)\le5$, on a
$ACS(h)\times I_{External}(h)\times M_{Controls}(h) \le 5\times1{,}2\times1{,}3=7{,}8$.
Par consequent, si $\mathrm{WeightedQDS}(h) < \tfrac{1000}{7{,}8}\approx128{,}2$,
alors $ARS(h)$ n'est \textbf{pas} sature, et $ARS(h)$ est exactement egal au produit
\cref{eq:ars} (le minimum n'intervient pas).
\end{propriete}

\begin{proof}
$M_{Controls}(h)=M_{SOC}(h)M_{AV}(h)M_{Patch}(h)$ est maximise lorsque le terme
protecteur $M_{SOC}$ et $M_{AV}$ valent chacun $1$ (absence de protection, valeur
neutre, superieure a $0{,}8$ et $0{,}9$) et que le terme aggravant $M_{Patch}$ vaut
$1{,}3$ (exclusion du cycle de patch) : $M_{Controls}(h)\le 1\times1\times1{,}3=1{,}3$.
Avec $ACS(h)\le5$ et $I_{External}(h)\le1{,}2$, le produit des trois facteurs est
majore par $5\times1{,}2\times1{,}3=7{,}8$. Si $\mathrm{WeightedQDS}(h)$ est tel que
$7{,}8\times\mathrm{WeightedQDS}(h) < 1000$, alors, a fortiori, le produit complet
\cref{eq:ars} (majore par $7{,}8\times \mathrm{WeightedQDS}(h)$) est strictement
inferieur a $1000$, donc $\min(\cdot,1000)$ vaut l'argument lui-meme.
\end{proof}

Une application numerique de cette propriete est fournie au \cref{chap:tests}
(\cref{sec:validation-numerique}).

\section{Analyse de sensibilite : la hierarchie des severites}

\begin{propriete}[Ratio de contribution entre severites]
\label{prop:ratio-severite}
A moyenne de $QDS$ et nombre de vulnerabilites egaux entre deux severites $i$ et $j$
($\mathrm{Avg}(QDS_i)=\mathrm{Avg}(QDS_j)$ et $n_i=n_j$), le rapport des termes
correspondants dans \cref{eq:qds-impl} est exactement egal au rapport des poids de
severite :
\[
\frac{\text{terme}_i}{\text{terme}_j} = \frac{w_i}{w_j}.
\]
Avec les valeurs de reference, une vulnerabilite Critique pese donc
\textbf{5 fois} plus qu'une vulnerabilite Elevee ($w_c/w_h=1{,}0/0{,}2=5$),
\textbf{20 fois} plus qu'une vulnerabilite Moyenne ($w_c/w_m=1{,}0/0{,}05=20$), et
\textbf{100 fois} plus qu'une vulnerabilite Faible ($w_c/w_l=1{,}0/0{,}01=100$).
\end{propriete}

\begin{proof}
Par definition, $\text{terme}_i = w_i\cdot\mathrm{Avg}(QDS_i)\cdot n_i^{1/100}$ et
$\text{terme}_j = w_j\cdot\mathrm{Avg}(QDS_j)\cdot n_j^{1/100}$. Sous l'hypothese
$\mathrm{Avg}(QDS_i)=\mathrm{Avg}(QDS_j)$ et $n_i=n_j$, ces deux facteurs communs se
simplifient dans le quotient, laissant $\text{terme}_i/\text{terme}_j = w_i/w_j$. Les
valeurs numeriques s'obtiennent en substituant $w_c=1{,}0$, $w_h=0{,}2$,
$w_m=0{,}05$, $w_l=0{,}01$.
\end{proof}

Cette propriete est la traduction mathematique exacte de l'intention metier du
modele : concentrer l'attention operationnelle sur les vulnerabilites critiques,
independamment de leur nombre relatif, puisque le rapport de poids est invariant et
ne depend jamais des donnees observees.

\section{Classification en classes de risque et gouvernance SLA}

\subsection{Classification}

L'ARS, une fois calcule, est projete sur cinq classes ordonnees par des seuils
configures $\tau_B>\tau_R>\tau_O>\tau_Y\ge0$ (valeurs de reference dans le
\cref{tab:sla}) :

\begin{equation}
\mathrm{Classe}(ARS) =
\begin{cases}
\text{BLACK} & \text{si } ARS \ge \tau_B \\
\text{RED} & \text{si } \tau_R \le ARS < \tau_B \\
\text{ORANGE} & \text{si } \tau_O \le ARS < \tau_R \\
\text{YELLOW} & \text{si } \tau_Y \le ARS < \tau_O \\
\text{GREEN} & \text{si } ARS < \tau_Y
\end{cases}
\label{eq:classify}
\end{equation}

\begin{table}[H]
\centering
\caption{Matrice dynamique des SLA de remediation (valeurs de reference)}
\label{tab:sla}
\begin{tabular}{@{}lllr@{}}
\toprule
\textbf{Classe} & \textbf{Seuil ARS} & \textbf{SLA exige} & \textbf{Couleur} \\
\midrule
Black  & $ARS \ge 700$          & 2 jours  & Noir \\
Red    & $600 \le ARS < 700$    & 7 jours  & Rouge \\
Orange & $400 \le ARS < 600$    & 15 jours & Orange \\
Yellow & $200 \le ARS < 400$    & 30 jours & Jaune \\
Green  & $ARS < 200$            & 90 jours & Vert \\
\bottomrule
\end{tabular}
\end{table}

\begin{propriete}[Bonne definition de la classification]
\label{prop:classify-total}
Si $\tau_B>\tau_R>\tau_O>\tau_Y\ge0$, alors \cref{eq:classify} definit une fonction
totale sur $[0,1000]$ : tout $ARS\in[0,1000]$ appartient a exactement une classe.
\end{propriete}

\begin{proof}
\emph{Exhaustivite.} Soit $ARS\in[0,1000]$. Si $ARS<\tau_Y$, il est classe GREEN. Sinon
$ARS\ge\tau_Y$ ; s'il est de plus $<\tau_O$, il est classe YELLOW. En poursuivant
ainsi, tout $ARS\ge\tau_Y$ non classe YELLOW/ORANGE/RED tombe necessairement dans le
dernier cas $ARS\ge\tau_B$ (BLACK), car les quatre cas couvrent respectivement
$(-\infty,\tau_Y)$, $[\tau_Y,\tau_O)$, $[\tau_O,\tau_R)$, $[\tau_R,\tau_B)$ et
$[\tau_B,+\infty)$, dont la reunion est $\mathbb{R}$ tout entier des lors que
$\tau_Y\le\tau_O\le\tau_R\le\tau_B$.
\emph{Exclusion mutuelle.} Ces cinq intervalles sont deux a deux disjoints car
chacun est de la forme $[\tau_{k},\tau_{k+1})$ avec les $\tau_k$ strictement
ordonnes (ou une demi-droite pour les cas extremes), donc un reel ne peut appartenir
qu'a un seul d'entre eux. Une valeur $ARS$ ne peut donc satisfaire qu'une seule
des cinq conditions de \cref{eq:classify}, ce qui prouve que la classification est
une fonction bien definie (et non une relation multivaluee).
\end{proof}

\begin{corollaire}[Coherence ordinale]
\label{cor:classify-ordre}
Si $ARS_1\le ARS_2$, la classe de $ARS_1$ n'est jamais strictement plus severe que
celle de $ARS_2$, pour l'ordre naturel $\text{GREEN}<\text{YELLOW}<\text{ORANGE}
<\text{RED}<\text{BLACK}$.
\end{corollaire}

\begin{proof}
Consequence immediate de \cref{eq:classify} : chaque classe correspond a un
intervalle de $ARS$, et ces intervalles sont ranges dans le meme ordre que la
severite des classes ($GREEN$ pour les plus petites valeurs, $BLACK$ pour les plus
grandes). La fonction $\mathrm{Classe}$ est donc une fonction en escalier
croissante de $ARS$ vers l'echelle ordinale des classes.
\end{proof}

C'est cette coherence qui garantit qu'augmenter le risque technique d'un actif ne
peut jamais faire baisser sa classe de gouvernance -- une propriete indispensable
pour la confiance des equipes dans le systeme.

\subsection{Gestion des exceptions (Risk Acceptance)}

Certains actifs ne peuvent etre patches dans les delais pour des raisons de
disponibilite metier ou d'incompatibilite logicielle. La CMDB porte alors un
attribut booleen \texttt{Exception}. Son activation \textbf{ne modifie pas} l'ARS
(le risque technique reste reel) mais suspend le decompte SLA standard :

\begin{equation}
\mathrm{StatutSLA}(h) =
\begin{cases}
\text{ACCEPTED\_RISK} & \text{si } \mathrm{Exception}(h) = \text{oui} \\
\text{VIOLATED} & \text{si non exceptee et delai depasse} \\
\text{ON\_TRACK} & \text{sinon}
\end{cases}
\label{eq:sla-status}
\end{equation}

ou $\text{date\_echeance}(h) = \text{premiere\_detection\_ouverte}(h) +
\mathrm{SLA}_{\mathrm{jours}}(\mathrm{Classe}(ARS(h)))$, le nombre de jours de SLA
etant lu dans le \cref{tab:sla} selon la classe de l'actif.

\section{Mesure de la performance : le MTTR pondere}

\subsection{Definition}

Pour un ensemble de vulnerabilites corrigees, regroupees par classe de risque $k$ de
l'actif auquel elles appartenaient, avec des poids de performance $w_k$
(configurables, valeurs de reference $w_{black}=10$, $w_{red}=5$, ...) :

\begin{equation}
\mathrm{MTTR}_w = \frac{\displaystyle\sum_k w_k \sum_{v \in \text{corrigees}_k} TTR(v)}
                       {\displaystyle\sum_k w_k \cdot C_k}
\label{eq:mttr}
\end{equation}

ou $C_k$ est le nombre de vulnerabilites corrigees de la classe $k$, et la somme ne
porte que sur les actifs \textbf{sans exception active} (\cref{eq:sla-status}),
conformement au principe de ne pas penaliser les equipes operationnelles pour des
blocages decides au niveau metier.

\begin{propriete}[Le MTTR pondere est une moyenne ponderee des MTTR par classe]
\label{prop:mttr-weighted-avg}
En notant $\mathrm{MTTR}_k = \frac{1}{C_k}\sum_{v\in\text{corrigees}_k} TTR(v)$ la
moyenne simple des temps de resolution de la classe $k$ (pour $C_k>0$), on a
\[
\mathrm{MTTR}_w = \sum_k \lambda_k\, \mathrm{MTTR}_k,
\qquad
\lambda_k = \frac{w_k C_k}{\sum_j w_j C_j} \ge 0,
\qquad
\sum_k \lambda_k = 1,
\]
et par consequent
\[
\min_k \mathrm{MTTR}_k \;\le\; \mathrm{MTTR}_w \;\le\; \max_k \mathrm{MTTR}_k .
\]
\end{propriete}

\begin{proof}
Par definition de $\mathrm{MTTR}_k$, $\sum_{v\in\text{corrigees}_k}TTR(v) = C_k\cdot
\mathrm{MTTR}_k$. En substituant dans le numerateur de \cref{eq:mttr} :
\[
\mathrm{MTTR}_w = \frac{\sum_k w_k\, C_k\, \mathrm{MTTR}_k}{\sum_k w_k C_k}
= \sum_k \underbrace{\frac{w_k C_k}{\sum_j w_j C_j}}_{\lambda_k}\, \mathrm{MTTR}_k .
\]
Chaque $\lambda_k\ge0$ car $w_k\ge0$ et $C_k\ge0$, et $\sum_k\lambda_k =
\frac{\sum_k w_k C_k}{\sum_j w_j C_j}=1$ : $(\lambda_k)_k$ est donc une distribution
de poids convexes, et $\mathrm{MTTR}_w$ est une combinaison convexe des
$\mathrm{MTTR}_k$. Toute combinaison convexe de reels est comprise entre leur
minimum et leur maximum : en notant $m=\min_k \mathrm{MTTR}_k$ et
$M=\max_k \mathrm{MTTR}_k$, on a $m\le \mathrm{MTTR}_k\le M$ pour tout $k$, donc
$m = \sum_k \lambda_k\, m \le \sum_k \lambda_k\, \mathrm{MTTR}_k \le \sum_k
\lambda_k\, M = M$.
\end{proof}

Cette propriete est essentielle pour l'interpretation de l'indicateur : le
$\mathrm{MTTR}_w$ affiche au comite de securite ne peut jamais etre une valeur
aberrante hors de la plage des performances reellement observees par classe ; il en
est une moyenne recentree sur les classes les plus critiques (celles qui recoivent
le plus grand poids $w_k$).

\section{Agregation globale : l'Enterprise Risk Score (ERS)}

\subsection{Definition}

Pour le pilotage strategique, les $ARS$ individuels sont agreges en un seul
indicateur, module par un facteur de performance operationnelle globale :

\begin{equation}
ERS = \left(\frac{1}{N}\sum_{i=1}^{N} ARS_i \right) \times \mathrm{Global\_I_{Ops}}
\label{eq:ers}
\end{equation}

\begin{equation}
\mathrm{Global\_I_{Ops}} =
\begin{cases}
1{,}00 & \text{si } \mathrm{MTTR}_{global} \le \mathrm{SLA}_{cible} \\
1{,}15 & \text{si } \mathrm{SLA}_{cible} < \mathrm{MTTR}_{global} \le 1{,}5\cdot \mathrm{SLA}_{cible} \\
1{,}30 & \text{si } \mathrm{MTTR}_{global} > 1{,}5\cdot \mathrm{SLA}_{cible}
\end{cases}
\label{eq:iops}
\end{equation}

\begin{propriete}[Monotonie du facteur de penalite]
\label{prop:iops-mono}
Pour $\mathrm{SLA}_{cible}>0$ fixe, $\mathrm{Global\_I_{Ops}}$ est une fonction
(faiblement) croissante de $\mathrm{MTTR}_{global}$.
\end{propriete}

\begin{proof}
Les trois cas de \cref{eq:iops} partitionnent $[0,+\infty)$ en trois intervalles
$\big(0,\mathrm{SLA}_{cible}\big]$, $\big(\mathrm{SLA}_{cible},1{,}5\,
\mathrm{SLA}_{cible}\big]$ et $\big(1{,}5\,\mathrm{SLA}_{cible},+\infty\big)$,
disjoints et ordonnes de gauche a droite. Sur chaque intervalle, la fonction est
constante, donc trivialement croissante. Aux deux points de raccord, la valeur
saute de $1{,}00$ a $1{,}15$ puis de $1{,}15$ a $1{,}30$ : la fonction ne decroit
jamais lorsqu'on passe d'un intervalle au suivant. Une fonction en escalier dont les
paliers sont croissants de gauche a droite est globalement croissante au sens large
sur son domaine.
\end{proof}

\begin{corollaire}[Monotonie de l'ERS en performance operationnelle]
Pour une moyenne d'ARS fixee, l'ERS est une fonction croissante de
$\mathrm{MTTR}_{global}$ : plus les equipes de remediation prennent du retard par
rapport a la cible SLA, plus l'indicateur global remonte a la direction est penalise,
independamment de toute amelioration ou degradation du risque technique brut.
\end{corollaire}

\begin{proof}
$ERS$ est le produit d'une constante positive ($\frac{1}{N}\sum ARS_i$, fixee par
hypothese) par $\mathrm{Global\_I_{Ops}}(\mathrm{MTTR}_{global})$, qui est croissante
par \cref{prop:iops-mono}. Le produit d'une constante strictement positive par une
fonction croissante est croissant.
\end{proof}

\begin{exemple}
Si la moyenne technique du parc est $\overline{ARS}=400$ (classe Orange) mais que le
$\mathrm{MTTR}_{global}$ mesure depasse $1{,}5$ fois la cible SLA (l'IT met
systematiquement deux fois plus de temps que prevu a corriger), alors
$\mathrm{Global\_I_{Ops}}=1{,}30$ et
\[
ERS = 400 \times 1{,}30 = 520,
\]
ce qui fait franchir a l'indicateur affiche a la direction la frontiere de la classe
Orange vers la classe Red (\cref{tab:sla}), reveletant une paralysie operationnelle
que la seule moyenne technique aurait masquee.
\end{exemple}

\subsection{Ecart entre specification et implementation actuelle}
\label{sec:ecart-ers}

A la date de redaction de ce rapport, le service \texttt{DashboardMetricsService}
calcule l'agregat expose au tableau de bord comme une simple moyenne arithmetique
des $ARS$ (\cref{eq:ers} sans le facteur $\mathrm{Global\_I_{Ops}}$) : aucune
cible SLA globale n'est encore configuree ni comparee au $\mathrm{MTTR}_w$ mesure
pour declencher la penalite \cref{eq:iops}. Le modele mathematique complet est donc
entierement specifie et demontre ci-dessus, mais son cablage dans le service de
calcul reste, a ce stade, une extension a realiser -- elle est reprise comme
perspective d'evolution prioritaire (\cref{sec:perspectives}).

\section{Synthese}

\begin{table}[H]
\centering
\caption{Recapitulatif des formules du modele de risque}
\label{tab:synthese-formules}
\begin{tabular}{@{}p{3.1cm}p{9.4cm}@{}}
\toprule
\textbf{Grandeur} & \textbf{Formule} \\
\midrule
$ACS(h)$ & $1 + 4\cdot\frac{\rho(\mathrm{role})+\varepsilon(\mathrm{env})}{2}$ \\[4pt]
$\mathrm{WeightedQDS}(h)$ & $\sum_i w_i\cdot\mathrm{Avg}(QDS_i)\cdot n_i^{1/100}$ \\[4pt]
$ARS(h)$ & $\min\big(ACS\times\mathrm{WeightedQDS}\times I_{External}\times M_{Controls}\,,\,1000\big)$ \\[4pt]
Classe$(ARS)$ & seuils $\tau_B>\tau_R>\tau_O>\tau_Y$ (\cref{eq:classify}) \\[4pt]
$\mathrm{MTTR}_w$ & $\dfrac{\sum_k w_k\sum TTR_k}{\sum_k w_k C_k}$ (hors actifs en exception) \\[4pt]
$ERS$ & $\Big(\frac1N\sum ARS_i\Big)\times \mathrm{Global\_I_{Ops}}(\mathrm{MTTR}_{global})$ \\
\bottomrule
\end{tabular}
\end{table}
